\section{Work Progress Overview}
\label{sec:work-progress}

The project developed through weekly milestones between 31 August and 11 November 2026. Table~\ref{tab:work-progress} summarises the deliverables reached at each stage, the main difficulties, and the corresponding adjustments. Khánh led dataset preparation and EDA; Tân and Chi developed the machine learning models; Phạm Thanh Bảo Ngân and Lê Phước Minh Thiện developed the neural models and prediction blends.

% Retrospective draft: weekly allocations and meeting topics are reconstructed
% from project deliverables, not verified against dated meeting records.
\begingroup
\small
\setlength{\emergencystretch}{3em}
\renewcommand{\arraystretch}{1.15}
\begin{longtable}{p{1.7cm}p{6.0cm}p{6.0cm}}
\caption[Weekly progress and difficulties]{Weekly progress, difficulties, and responses.}
\label{tab:work-progress}\\
\toprule
Week & Progress reached & Difficulty and response \\
\midrule
\endfirsthead
\multicolumn{3}{l}{\textit{Table~\ref{tab:work-progress} continued.}}\\
\toprule
Week & Progress reached & Difficulty and response \\
\midrule
\endhead
\midrule
\multicolumn{3}{r}{\textit{Continued on the next page.}}\\
\endfoot
\bottomrule
\endlastfoot

\textbf{1}\newline
31 Aug.--6 Sep. & Agreed on asking-price prediction, selected the Vietnamese real estate dataset, reviewed the available attributes, and divided the work among five members. & Listings include both sale and rental advertisements without a reliable transaction-type label. The group defined the target as advertised asking price and retained this ambiguity as a limitation. \\
\midrule

\textbf{2}\newline
7--13 Sep. & Completed initial loading, schema checks, data-type inspection, and missing-value summaries. Established a reusable local data copy for subsequent analysis. & The dataset contains approximately 3.5 million listings, making repeated loading costly. Local Parquet storage and categorical conversion reduced repeated data-access and memory overhead. \\
\midrule

\textbf{3}\newline
14--20 Sep. & Completed duplicate detection and checks for invalid price and area values. Produced missing-value summaries by property type. & Many structural attributes are absent, and missingness differs by property type. The group avoided dropping every incomplete listing and retained missingness for model-specific treatment. \\
\midrule

\textbf{4}\newline
21--27 Sep. & Completed the main price and area analyses, reviewed extreme observations, and established the cleaned modelling population of 2,397,001 listings. & Extreme prices, implausible dimensions, and invalid timestamps affected modelling suitability. Price, area, and price-per-area restrictions were applied; invalid optional dimensions were set to missing. \\
\midrule

\textbf{5}\newline
28 Sep.--4 Oct. & Completed hierarchical location keys, categorical missing-value handling, text preparation, and the shared chronological 70/15/15 split. & Repeated district names could conflate locations, while advertisements could disclose the target price. Parent-location prefixes and monetary-clause/digit removal addressed these issues, with residual text leakage retained as a limitation. \\
\midrule

\textbf{6}\newline
5--11 Oct. & Established constant-median and local price-per-square-metre baselines. Completed structured Ridge training and the three-value regularisation search. & Linear fitting required consistent treatment of skewed numeric inputs, missing values, and numerous categorical levels. Training-fitted log transforms, imputation, scaling, and rare-category pooling formed a common pipeline. \\
\midrule

\textbf{7}\newline
12--18 Oct. & Completed the Random Forest depth/feature-fraction comparison and the LightGBM, CatBoost, and XGBoost training configurations. Prepared validation metrics and feature-importance diagnostics. & Tree experiments on the full training population were computationally demanding. A 10\% bootstrap fraction per forest tree, parallel CPU execution, and validation-based stopping controls limited the training workload. \\
\midrule

\textbf{8}\newline
19--25 Oct. & Completed TF-IDF features, joint text-and-structured Ridge, and the text-prediction input for LightGBM. Completed the tabular MLP architecture and its preprocessing pipeline. & Dense text matrices would be costly, and in-sample text predictions could mislead the downstream model. Sparse matrices and five-fold out-of-fold Ridge predictions were used; categorical embeddings reduced the neural input expansion. \\
\midrule

\textbf{9}\newline
26 Oct.--1 Nov. & Developed the MLP optimisation and early-stopping procedure. Prepared BamiBERT tokenisation, standardised log-price targets, and GPU fine-tuning. & Neural training required control of sequence length, numerical scale, and GPU workload. Text was limited to 512 tokens, targets were standardised for BamiBERT, and mixed precision was configured. \\
\midrule

\textbf{10}\newline
2--8 Nov. & Consolidated neural-model predictions, compared structured and text-based approaches, and developed nonnegative prediction blends with validation-fitted coefficients. & Combining separately trained models required aligned observations and consistent prediction units. Validation/test predictions retained common row identities and log-price units. Separate reconstruction of text-model ordering remained a consistency check. \\
\midrule

\textbf{11}\newline
9--11 Nov. & Brought together experiment summaries, diagnostic figures, and report contributions. Reviewed model settings, interpretation, and consistency between individual and combined predictions. & Results from different experimental populations could not be pooled directly. Comparisons were restricted to compatible runs; remaining text-filtering and checkpoint-compatibility limitations were retained for discussion. \\

\end{longtable}
\endgroup

\paragraph{Fortnightly tracking and discussion.}
Following the kickoff on 31 August, the tracking milestones were set two weeks apart: 14 September (data loading and quality issues), 28 September (cleaning decisions and feature preparation), 12 October (baselines and evaluation protocol), 26 October (tree/text results and neural-training requirements), and 9 November (prediction alignment, model combinations, and report integration). At each checkpoint, members reviewed deliverables, discussed unresolved issues, and agreed on the next tasks. A final review on 11 November covered the consolidated findings and outstanding limitations. Intermediate outputs were also exchanged between meetings so that modelling work could proceed in parallel.
